\section{Mechanistic analysis and result logic}

For every MoE we measure expert usage, assignment entropy, load imbalance, and dead experts. On a
held-out in-distribution audit set, we estimate acquisition-environment and biological-label
information in routes using mutual information and cross-validated decoders. We apply identical
linear probes to final embeddings to measure environment and label decodability.

We also replace learned assignments with seeded random assignments while freezing expert weights.
The accuracy change measures route reliance: a model can exhibit structured routing without
requiring that routing for its predictions. We relate conditional gain to route reliance,
environment decodability, label decodability, gradient conflict, and utilization across held-out
experiments. These are descriptive mechanism associations; causal claims are restricted to the
matched architectural contrast.

A positive matched gain with high route reliance supports useful conditional specialization. A
gain only against the smaller original encoder indicates capacity, not conditionality. High
acquisition decodability without matched gain indicates nuisance partitioning. If gradients
conflict but matched conditional gain remains absent across the held-out-experiment folds, then
sparse routing does not solve the interference mechanism proposed here.

\paragraph{Substrate qualification.}
Table~\ref{tab:native-competence} reports the predeclared native-stain Cell-DINO qualification.
Full fine-tuning recovers substantial task signal relative to a frozen linear readout, but its
performance drops by 14.80 points from seen experiments to OOD validation. The substrate is thus
competent enough to test conditional capacity, and the intended failure regime is batch transfer
rather than a floor representation. This result does not establish the proposed mechanism: the
same gap could arise from ordinary overfitting or acquisition variation that routing cannot fix.

\begin{table}[t]
  \centering
  \caption{Native-stain Cell-DINO qualification on RxRx1 (\%). OOD test is untouched.}
  \label{tab:native-competence}
  \begin{tabular}{lrrrr}
    \toprule
    Adaptation & Train & ID & OOD-val & Worst experiment \\
    \midrule
    Frozen backbone + linear head & 10.62 & 7.64 & 4.42 & 0.81 \\
    Full fine-tuning & 39.00 & 27.09 & 12.29 & 1.42 \\
    \bottomrule
  \end{tabular}
\end{table}

Table~\ref{tab:kill-contrast} reports the predeclared seed-0 kill contrast. Relative to the
exact-total-parameter-matched dense-wide model, MoE improves OOD-validation accuracy by 1.15 points,
ID accuracy by 1.98 points, and worst-experiment accuracy by 0.16 points. All four held-out
experiments move in the same direction. The gain is nevertheless far below the frozen 5-point
replication trigger, and the smaller original encoder remains 0.70 OOD-validation points better
than MoE. We therefore stop replication and architecture expansion. This single-seed gate does not
establish a zero population effect; it rejects the targeted 10--15-point effect strongly enough to
stop the predeclared campaign.

\begin{table}[t]
  \centering
  \caption{Seed-0 native-stain RxRx1 kill contrast (accuracy in \%). Dense-wide and MoE differ by
  0.001232\% in total parameters. OOD test is untouched.}
  \label{tab:kill-contrast}
  \begin{tabular}{lrrrrrr}
    \toprule
    Model & Params (M) & Active FFN (M) & Train & ID & OOD-val & Worst \\
    \midrule
    Original & 22.40 & 1.182 & 39.12 & 26.99 & 12.31 & 1.50 \\
    Dense-wide & 30.68 & 9.455 & 36.13 & 24.36 & 10.46 & 1.22 \\
    MoE & 30.68 & 1.185 & 38.72 & 26.35 & 11.61 & 1.38 \\
    \bottomrule
  \end{tabular}
\end{table}

The result does not support robustness from conditional capacity beyond ordinary capacity: MoE
recovers part of the degradation caused by dense widening but does not surpass the original
substrate.

The clean bounded routing diagnosis in Table~\ref{tab:route-diagnosis} further disfavors H2. The
learned router uses all eight experts, but its normalized usage entropy is essentially maximal and
route mutual information with experiment and perturbation is near zero. Randomizing its routes
reduces OOD-validation accuracy by only 0.42 points. A router frozen at initialization is 0.19
points better than the learned router and remains less than one point route-dependent. Learning the
router therefore does not identify a reusable batch-specialized partition that is necessary for
the model's predictions.

\begin{table}[t]
  \centering
  \scriptsize
  \caption{Clean seed-0 route diagnosis on OOD validation (accuracy and reliance in percentage
  points). Both controls use all eight experts. MI denotes the aligned token-route audit.}
  \label{tab:route-diagnosis}
  \begin{tabular}{lrrrrrr}
    \toprule
    Router & OOD-val & Random routes & Reliance & Entropy & MI(exp.) & MI(label) \\
    \midrule
    Learned & 11.65 & 11.23 & 0.42 & 0.9999 & 0.0051 & 0.0016 \\
    Frozen & 11.84 & 11.13 & 0.71 & 0.9975 & 0.0098 & 0.0041 \\
    \bottomrule
  \end{tabular}
\end{table}

This mechanism result is diagnostic rather than a population efficacy estimate. It is compatible
with sparse activation weakly regularizing a harmful dense widening, a fixed random partition
creating mildly useful subspaces, or single-seed variation. It does not establish the proposed
cross-batch gradient-interference mechanism, and it cannot reopen the failed replication gate. All
OOD-test values remain withheld.

\paragraph{Long-adaptation substrate-strength test.}
The bounded 90-epoch matrix in Table~\ref{tab:hypothesis90} separates a weak substrate from a weak
sparse intervention. Every adapted arm fits the training set perfectly and reaches 48--53\% ID
accuracy, while OOD-validation accuracy remains 18.5--20.5\%. Cell-DINO is therefore a competent
instrument and the remaining failure is transfer to unseen experiments, not inability to learn the
training task.

Canonical top-1 MoE improves OOD validation by 1.72 points and ID by 3.21 points relative to the
matched dense-wide model, but worst-experiment accuracy falls by 0.08 points. It is only 0.12 OOD
points above the smaller original model, which has a better worst-experiment score. The matched
models differ by 378 parameters (0.001232\%), within the predeclared 0.1\% tolerance. The primary
contrast therefore fails the five-point replication gate.

The matched OOD gain is $-0.42$, $+2.78$, $+2.55$, and $+1.72$ points at epochs 10, 30, 60, and
90, respectively. The conclusion is therefore already visible by epoch 60 and does not arise from
stopping before a growing sparse advantage; the final 30 epochs strengthen absolute baselines
while the conditional gain contracts.

\begin{table}[t]
  \centering
  \caption{Validated seed-0 90-epoch RxRx1 hypothesis matrix (accuracy in \%). Top-2 activates
  two experts and is not active-compute matched; partial and loss-based arms are diagnostic. OOD
  test is untouched.}
  \label{tab:hypothesis90}
  \begin{tabular}{lrrrr}
    \toprule
    Arm & Train & ID & OOD-val & Worst \\
    \midrule
    Original anchor & 100.00 & 52.09 & 20.09 & 1.75 \\
    Matched dense-wide & 100.00 & 48.23 & 18.50 & 1.70 \\
    Canonical token top-1 MoE & 100.00 & 51.45 & 20.22 & 1.62 \\
    Image top-1 MoE & 100.00 & 50.45 & 20.13 & 1.46 \\
    Within-experiment-balanced MoE & 100.00 & 51.96 & 20.46 & 1.38 \\
    Token top-2 MoE & 100.00 & 49.91 & 19.79 & 1.50 \\
    Frozen linear & 31.77 & 14.73 & 6.69 & 0.53 \\
    Last-four-block adaptation & 100.00 & 53.01 & 19.82 & 1.79 \\
    Output invariance & 100.00 & 43.53 & 16.58 & 1.38 \\
    Environment-balanced classification & 100.00 & 50.85 & 19.99 & 1.66 \\
    \bottomrule
  \end{tabular}
\end{table}

The alternative arms do not expose a hidden large sparse effect. Top-2 trails top-1, weakening
hard-routing starvation as the dominant explanation. Image routing is effectively tied with token
routing, so token granularity is not decisive. Within-experiment balancing gives the highest mean
OOD accuracy but the lowest worst-experiment score among the main adapted arms. Last-four-block
adaptation has the best ID and worst-experiment values without the best mean OOD value. These
patterns retain regularization, representation depth, and single-seed variation as alternatives;
they do not support a uniquely sparse or tail-robust mechanism. No replication, deeper mechanism
campaign, or OOD-test evaluation is opened by this result.

\paragraph{Exploratory post-gate confirmation.}
Under a subsequent explicitly multiplicity-exposed architecture search, the tail-safe seed-0
recipe was locked and compared with its exact-total-parameter-matched dense control at fresh seeds
1 and 2. At epoch 30, sparse minus dense changes OOD-validation/ID/worst-experiment accuracy by
$+1.76/+1.89/+0.04$ points for seed 1 and $+1.45/+1.50/+0.12$ points for seed 2. Thus the signs
align across both fresh seeds and all three decision axes, but the mean OOD effect remains far
below five points. These milestones are not a new tuned winner or a completed population claim:
all four locked arms continue unchanged to epoch 60, and the OOD test remains untouched. The
principal threat is selection across the preceding seed-0 factorial, router-auxiliary, and
temperature screens; contraction or sign reversal at epoch 60 is the predeclared falsifier.
